
% natural queries -> the semi-structure knowledge bases
% the semi-structure knowledge bases -> retrieval
% retrieval -> benchmark retrieval
% Why in the world of llms we still need this benchmark 
% Why challenging -> need to infer and reaon and combine - before 'finally'
% Dif problem motivate better.
% \hw{Note: Submission is single-blind. Contents about human queries are highlighted.}

Answering real-world complex queries, such as complex product search, often requires accurate retrieval from semi-structured knowledge bases that involve blend of unstructured (\eg textual descriptions of products) and structured (\eg entity relations of products) information. 
% Real-world queries often feature a sophisticated blend of unstructured (\eg textual descriptions) and structured (\eg entity relations \qian{database of products?}) information, which naturally involve semi-structured knowledge bases from specific domains. This highlights the essential need for retrieving knowledge from these semi-structured knowledge bases to accurately answer user queries.
% \michi{We could be more clear and crisp. how about something like: "Answering real-world user queries, such as product search, often requires accurate retrieval of information from semi-structured knowledge bases or databases that involve
% blend of unstructured (\eg textual descriptions of products) and structured (\eg entity relations of products) information"}
However, many previous works studied textual and relational retrieval tasks as separate topics.
To address the gap, we develop \benchmarkt, a large-scale \underline{S}emi-structure retrieval benchmark on \underline{T}extual \underline{a}nd \underline{R}elational \underline{K}nowledge Bases.
Our benchmark covers three domains: product search, academic paper search, and queries in precision medicine.
We design a novel pipeline to synthesize realistic user queries that integrate diverse relational information and complex textual properties, together with their ground-truth answers (items). 
We conduct rigorous human evaluation to validate the quality of our synthesized queries. 
\hw{We further enhance the benchmark with high-quality human-generated queries to provide an authentic reference.}
\benchmarkh serves as a comprehensive testbed for evaluating the performance of retrieval systems driven by large language models (LLMs). Our experiments suggest that \benchmarkh presents significant challenges to the current retrieval and LLM systems, 
%\michi{you mentioned LLM earlier in your abstract, so we should revisit it here}
highlighting the need for more capable semi-structured retrieval systems. %  that can handle both textual and relational aspects
% The benchmark data and code are available on \textcolor{gray}{\textcolor{cyan}{\url{https://github.com/snap-stanford/stark}}}.

% \jure{can we post URLs? isnt' it double blind?}\shirley{The submission is single-blind by default "Reviewing is in principle single-blind, hence the paper should not be anonymized. In cases where the work can be reviewed equally well anonymously, anonymous submission is also allowed." \url{https://neurips.cc/Conferences/2024/CallForDatasetsBenchmarks}. Please let me know if we want to do double-blind instead}